% !TeX root = ../../../main.tex
\section{Tame Ramification Along the Full Boundary}\label{section:tame-ramification-full-boundary}

The ramification calculation of Chapter~2 proves the required tameness when
the modulus is supported at a single point and the degree is the degree of
that modulus. We now remove these two restrictions. First, we pass from the
point-supported components of $\mdls$ to the full modulus in degree
$\deg\mdls$. We then compare the blowup at the point
$\mdls\in C^{(\deg\mdls)}$ with the blowup in degree $d$ along the locus of
divisors containing $\mdls$. This proves the following result.

\begin{theorem}\label{theorem:SymmetricPowerOfSheafIsTamelyRamified}
Let $G$ be a finite abelian group, and let $\cP$ be a $G$-torsor on $U$
whose ramification is bounded by $\mdls$. For every sufficiently large
integer $d$, the torsor $\cP^{[d]}$ is tamely ramified along
\[
H=\Cmod{d}\setminus U^{(d)}
  =E_0\times_{X_{\mdls,d}}\Cmod{d}
  \cong Z_0\times_{C^{(d)}}\Cmod{d}.
\]
\end{theorem}

To prove the theorem, we retain the blowup notation of Chapter~2. For every
nonzero effective submodulus $\ndls\leq\mdls$, we identify $\ndls$ with the
corresponding $k$-point of $C^{(\deg\ndls)}$, and put
\[
    X_{\ndls}=\Bl_{\ndls}\bigl(C^{(\deg\ndls)}\bigr),
    \qquad
    E_{\ndls}=X_{\ndls}\times_{C^{(\deg\ndls)}}\{\ndls\}.
\]
We write $\eta_{\ndls}$ for the generic point of $E_{\ndls}$. For the full
modulus in degree $d$, the center is
\[
Z_0=Z_0(\mdls,d)=\mdls+C^{(d-\deg\mdls)}\subset C^{(d)},
\]
where the right-hand side denotes the image of the closed immersion
$D\mapsto\mdls+D$.
The corresponding blowup, exceptional divisor, and generic point are denoted
by $X_{\mdls,d}$, $E_0=E_{\mdls,d}$, and
$\eta_0=\eta_{\mdls,d}$, as in
\Cref{section:BlowupOfSymmetricPowerOfCurves}.

\noindent\textbf{From point-supported moduli to the full modulus.}

We first show that if the relevant torsors are tame for two coprime
submoduli, then the corresponding torsor is tame for their sum. Iterating
this step over the point-supported components of $\mdls$ yields the required
statement for the full modulus in degree $\deg\mdls$.

To establish this implication, we use the following general comparison. Let
$X$ and $Y$ be smooth integral schemes over a field $k$, and let
$x\in X(k)$ and $y\in Y(k)$.
We denote the blowups of these schemes at the given points by
$\pi_X: \Bl_x(X) \to X$ and $\pi_Y: \Bl_y(Y) \to Y$. Furthermore, let $\pi_{X \times Y}: \Bl_{(x,y)}(X \times_k Y) \to X \times_k Y$
be the blowup of the product scheme at the point $(x,y)$.
We denote by $E_X, E_Y$, and $E_{X \times Y}$ the respective exceptional divisors, and let $\eta_X, \eta_Y$, and $\eta_{X \times Y}$ be their generic points.

Tame ramification is preserved under the contracted product of torsors after
pullback to a product, as follows.

\begin{prop}\label{prop:ramification-external-product}
Let $G$ be a finite abelian group. Suppose $\cP_X$ and $\cP_Y$ are
$G$-torsors defined on dense open subsets
$U_X\subset\Bl_x(X)$ and $U_Y\subset\Bl_y(Y)$ disjoint from the respective
exceptional divisors. There exist a dense open subset
\[
    \widetilde U\subset\Bl_{(x,y)}(X\times_kY)
\]
containing $\eta_{X\times Y}$ and morphisms
\[
    r_X:\widetilde U\longrightarrow\Bl_x(X),
    \qquad
    r_Y:\widetilde U\longrightarrow\Bl_y(Y),
\]
which send $\eta_{X\times Y}$ to $\eta_X$ and $\eta_Y$, respectively. On
the dense open subset
\[
    W=r_X^{-1}(U_X)\cap r_Y^{-1}(U_Y)\subset\widetilde U,
\]
define
\[
    \cP_{X\times Y}
      :=r_X^{-1}\cP_X
        \otimes^G
        r_Y^{-1}\cP_Y.
\]
If $\cP_X$ is tamely ramified at $\eta_X$ and $\cP_Y$ is tamely ramified
at $\eta_Y$, then $\cP_{X\times Y}$ is tamely ramified at
$\eta_{X\times Y}$.
\end{prop}

\begin{proof}
We first construct $r_X$ and $r_Y$; the torsors enter only after these
morphisms have been defined on the original blowups. Shrinking $X$ and $Y$
around $x$ and $y$ does not change the exceptional divisors or their local
rings at the generic points. The \`etale local description of smooth
morphisms \stackstag{054L} therefore gives \`etale morphisms
\[
a:X\longrightarrow\mathbb A_k^n,
\qquad
b:Y\longrightarrow\mathbb A_k^m
\]
carrying $x$ and $y$ to the origins; after shrinking further, we may assume
that $x$ and $y$ are the scheme-theoretic inverse images of the respective
origins. We write $x_1,\ldots,x_n$ and $y_1,\ldots,y_m$ for the corresponding
standard coordinates and also for their pullbacks to $X$ and $Y$.

Write
\[
    B_n=\operatorname{Bl}_0(\mathbb A_k^n),\qquad
    B_m=\operatorname{Bl}_0(\mathbb A_k^m),\qquad
    B_{n,m}=\operatorname{Bl}_{(0,0)}
        (\mathbb A_k^n\times_k\mathbb A_k^m),
\]
and denote their exceptional divisors by $E_n$, $E_m$, and $E_{n,m}$.
The blowup $B_{n,m}$ is the subscheme of
$(\mathbb A_k^n\times_k\mathbb A_k^m)\times_k\mathbb P_k^{n+m-1}$ defined by

\[\begin{cases}
x_i w_j = x_j w_i & 1 \le i, j \le n \\
y_k w_{n+l} = y_l w_{n+k} & 1 \le k, l \le m \\
x_i w_{n+j} = y_j w_i & 1 \le i \le n, \,\, 1 \le j \le m
\end{cases}\]
where $[w_1:\dots:w_{n+m}]$ are the homogeneous coordinates of
$\mathbb P_k^{n+m-1}$. The exceptional divisor $E_{n,m}$ is
isomorphic to $\mathbb P_k^{n+m-1}$.

Both $B_n\times_kB_m$ and $B_{n,m}$ are birational to
$\mathbb A_k^n\times_k\mathbb A_k^m$; their common dense open is
\[
    (\mathbb A_k^n\setminus\{0\})
      \times_k(\mathbb A_k^m\setminus\{0\}).
\]
Let
\[
    q_{n,m}\colon B_{n,m}
      \longrightarrow\mathbb P_k^{n+m-1}
\]
be the second projection. Define
\[
    \Omega_n=\bigcup_{i=1}^{n}D_+(w_i),
    \qquad
    \Omega_m=\bigcup_{j=1}^{m}D_+(w_{n+j}),
\]
and
\[
    \widetilde U_{n,m}
      :=q_{n,m}^{-1}(\Omega_n\cap\Omega_m)
      =\bigcup_{\substack{1\leq i\leq n\\1\leq j\leq m}}
        q_{n,m}^{-1}\!\left(D_+(w_iw_{n+j})\right).
\]
There is a morphism
\[
\begin{aligned}
f_{n,m}\colon\widetilde U_{n,m}&\longrightarrow
   B_n\times_kB_m,\\
((u,v),[w_1:\cdots:w_{n+m}])&\longmapsto
\bigl((u,[w_1:\cdots:w_n]),
      (v,[w_{n+1}:\cdots:w_{n+m}])\bigr).
\end{aligned}
\]
Thus we have a diagram
\[
\begin{tikzcd}
& \widetilde U_{n,m}
    \arrow[dl,"f_{n,m}"'] \arrow[dr,hook] & \\
B_n\times_kB_m && B_{n,m}.
\end{tikzcd}
\]
The restriction of $q_{n,m}$ induces an isomorphism
$E_{n,m}\xrightarrow{\sim}\mathbb P_k^{n+m-1}$. Since neither
$V_+(w_1,\ldots,w_n)$ nor $V_+(w_{n+1},\ldots,w_{n+m})$ contains the generic
point of $\mathbb P_k^{n+m-1}$, we have
$\eta_{n,m}\in\widetilde U_{n,m}$, where $\eta_{n,m}$ is the generic point
of $E_{n,m}$.
On the exceptional divisors, the restriction of $f_{n,m}$ is the
coordinate-block morphism
\[
\begin{aligned}
E_{n,m}\cap\widetilde U_{n,m}&\longrightarrow E_n\times_kE_m,\\
[w_1:\cdots:w_{n+m}]&\longmapsto
\bigl([w_1:\cdots:w_n],[w_{n+1}:\cdots:w_{n+m}]\bigr).
\end{aligned}
\]
This morphism is surjective and hence dominant. It therefore sends
$\eta_{n,m}$ to the generic point $(\eta_n,\eta_m)$ of $E_n\times_kE_m$.

We now return to the original schemes. By
\Cref{theorem:blowup_flat_base_change}, there are canonical identifications
\[
\begin{aligned}
\operatorname{Bl}_x(X)&\cong X\times_{\mathbb A_k^n} B_n,\\
\operatorname{Bl}_y(Y)&\cong Y\times_{\mathbb A_k^m} B_m,\\
B:=\operatorname{Bl}_{(x,y)}(X\times_kY)
  &\cong (X\times_kY)
      \times_{\mathbb A_k^n\times_k\mathbb A_k^m}B_{n,m}.
\end{aligned}
\]
Set
\[
    \widetilde U
      :=(X\times_kY)
        \times_{\mathbb A_k^n\times_k\mathbb A_k^m}\widetilde U_{n,m}.
\]
It is a dense open neighborhood of $\eta_{X\times Y}$ in $B$. Base change of
$f_{n,m}$ gives a morphism
\[
    f:\widetilde U\longrightarrow
      \operatorname{Bl}_x(X)\times_k\operatorname{Bl}_y(Y)
\]
that sends $\eta_{X\times Y}$ to $(\eta_X,\eta_Y)$. If $p_X$ and $p_Y$ are
the two projections from its target, set
\[
    r_X=p_X\circ f,
    \qquad
    r_Y=p_Y\circ f.
\]
These are the morphisms asserted in the statement. They fit into the
commutative diagram
\[
\begin{tikzcd}[column sep=large,row sep=large]
& \operatorname{Bl}_{(x,y)}(X\times_kY) & \\
& \widetilde U
    \arrow[u,hook] \arrow[dl,"r_X"'] \arrow[d,"f"] \arrow[dr,"r_Y"] & \\
\operatorname{Bl}_x(X)
& \operatorname{Bl}_x(X)\times_k\operatorname{Bl}_y(Y)
    \arrow[l,"p_X"] \arrow[r,"p_Y"']
& \operatorname{Bl}_y(Y).
\end{tikzcd}
\]
Over the common dense
open $(X\setminus\{x\})\times_k(Y\setminus\{y\})$, the morphism $f$ is the
evident identification; in particular, it is dominant. Hence
\[
    W=r_X^{-1}(U_X)\cap r_Y^{-1}(U_Y)
      =f^{-1}(U_X\times_k U_Y)
\]
is a nonempty open subset of the integral scheme $\widetilde U$, and is
therefore dense in $\widetilde U$ and in $B$.

It remains to compare the maps induced by $r_X$ and $r_Y$ on the exceptional
valuations. Let
\[
S=\mathcal O_{B,\eta_{X\times Y}},\qquad
R_X=\mathcal O_{\operatorname{Bl}_x(X),\eta_X},\qquad
R_Y=\mathcal O_{\operatorname{Bl}_y(Y),\eta_Y}.
\]
The morphisms $r_X$ and $r_Y$ induce local homomorphisms
\[
    R_X\longrightarrow S,
    \qquad
    R_Y\longrightarrow S.
\]
By base change, the exceptional divisors are identified with
\[
    E_X\cong\mathbb P_k^{n-1},\qquad
    E_Y\cong\mathbb P_k^{m-1},\qquad
    E_{X\times Y}\cong\mathbb P_k^{n+m-1},
\]
and the restrictions of $r_X$ and $r_Y$ are the two coordinate projections
\[
[w_1:\cdots:w_{n+m}]\longmapsto[w_1:\cdots:w_n],
\qquad
[w_1:\cdots:w_{n+m}]\longmapsto[w_{n+1}:\cdots:w_{n+m}].
\]
At the generic point of $E_{X\times Y}$ both coordinate blocks are nonzero.
On the chart $D_+(w_1w_{n+1})$, the functions $x_1$ and $y_1$ are
uniformizers in $R_X$ and $R_Y$, respectively, and each is also a
uniformizer in $S$. Hence both extensions have ramification index one. Their
residue-field extensions are
\[
k\left(\frac{x_2}{x_1},\ldots,\frac{x_n}{x_1}\right)
\subset
k\left(\frac{x_2}{x_1},\ldots,\frac{x_n}{x_1},
        \frac{y_1}{x_1},\ldots,\frac{y_m}{x_1}\right)
\]
and
\[
k\left(\frac{y_2}{y_1},\ldots,\frac{y_m}{y_1}\right)
\subset
k\left(\frac{x_1}{y_1},\ldots,\frac{x_n}{y_1},
        \frac{y_2}{y_1},\ldots,\frac{y_m}{y_1}\right),
\]
respectively. In particular, both are purely transcendental and hence
separable.

We now apply this comparison to the torsors. The extensions
$R_X\to S$ and $R_Y\to S$ are weakly unramified and have
separable residue-field extensions. Hence
\Cref{lemma:weakly_unramified_base_change} shows that the pullbacks of
$\cP_X$ and $\cP_Y$ to $\operatorname{Frac}(S)$ are tamely ramified.
By the definition above, their contracted product is the generic fiber of
$\cP_{X\times Y}$. Therefore
\Cref{lemma:bounding_ramification_of_contracted_product} shows that
$\cP_{X\times Y}$ is tamely ramified at $\eta_{X\times Y}$.
\end{proof}

\begin{cor}\label{cor:moduli-reduction}
Let $G$ be a finite abelian group, and let $\cP$ be a $G$-torsor on $U$.
Let $\ndls_1,\ndls_2\subset\mdls$ be two coprime nonzero submoduli.
Assume that $\cP^{[\deg\ndls_1]}$ and $\cP^{[\deg\ndls_2]}$ are tamely
ramified at $\eta_{\ndls_1}$ and $\eta_{\ndls_2}$, respectively. Then
$\cP^{[\deg\ndls_1+\deg\ndls_2]}$ is tamely ramified at
$\eta_{\ndls_1+\ndls_2}$.
\end{cor}

\begin{proof}
Put $d_i=\deg\ndls_i$ for $i=1,2$, and pull $\cP^{[d_i]}$ back to
$X_{\ndls_i}$ over the inverse image of $U^{(d_i)}$. Applying
\Cref{prop:ramification-external-product} to the two blowups
$X_{\ndls_i}\to C^{(d_i)}$, and writing $r_i$ for the morphisms supplied
by that proposition, shows that
\[
r_1^{-1}\cP^{[d_1]}\otimes^G r_2^{-1}\cP^{[d_2]}
\]
is tame; i.e., this contracted product has tame ramification at the generic
point of the exceptional divisor of
\[
\Bl_{(\ndls_1,\ndls_2)}
   \bigl(C^{(d_1)}\times_k C^{(d_2)}\bigr).
\]

By \Cref{lemma:takeuchi_lemma}, applied with
$\mdls=\ndls_1+\ndls_2$ and $d=d_1+d_2$, the addition morphism
\[
+_{d_1,d_2}\colon
C^{(d_1)}\times_k C^{(d_2)}
\longrightarrow C^{(d_1+d_2)}
\]
is \`etale at $(\ndls_1,\ndls_2)$ and maps this point to
$\ndls_1+\ndls_2$. Choose open neighborhoods
\[
(\ndls_1,\ndls_2)\in V\subset C^{(d_1)}\times_k C^{(d_2)},
\qquad
\ndls_1+\ndls_2\in W\subset C^{(d_1+d_2)}
\]
such that $+_{d_1,d_2}|_V\colon V\to W$ is \`etale and
\[
V\times_W\{\ndls_1+\ndls_2\}=\{(\ndls_1,\ndls_2)\}
\]
scheme-theoretically. Flat base change for blowups gives
\[
\Bl_{(\ndls_1,\ndls_2)}(V)
\cong
V\times_W\Bl_{\ndls_1+\ndls_2}(W).
\]
Hence the second projection defines an \`etale morphism
\[
\widetilde{+}_{d_1,d_2}\colon
\Bl_{(\ndls_1,\ndls_2)}(V)
\longrightarrow
\Bl_{\ndls_1+\ndls_2}(W)
\subset X_{\ndls_1+\ndls_2}.
\]
If $\eta$ denotes the generic point of the exceptional divisor of the
source, then
\[
\widetilde{+}_{d_1,d_2}(\eta)=\eta_{\ndls_1+\ndls_2}.
\]
By \Cref{prop:symmetric_powers_on_curves}, on the common dense open,
\[
(\widetilde{+}_{d_1,d_2})^{-1}\cP^{[d_1+d_2]}
\cong
r_1^{-1}\cP^{[d_1]}\otimes^G r_2^{-1}\cP^{[d_2]}.
\]
Its tameness therefore descends by
\Cref{lemma:descend_ramification_along_etale}.
\end{proof}

\noindent\textbf{From the point center to the degree-$d$ center.}

The comparison is governed by the addition morphism
\[
+\colon C^{(\deg\mdls)}\times_k C^{(d-\deg\mdls)}\longrightarrow C^{(d)}.
\]
After restricting the second factor to $U^{(d-\deg\mdls)}$ and blowing up,
it induces a comparison morphism
\[
X_{\mdls}\times_k U^{(d-\deg\mdls)}\longrightarrow X_{\mdls,d}.
\]
To descend tameness along this morphism, we must show that it is \`etale at
the generic point of the exceptional divisor in its source. This follows
from the case $l=1$ of the next lemma, adapted from
\cite[Lemma~4.1]{takeuchi2018blow}.

\begin{lemma}\label{lemma:takeuchi_lemma}
     Let $C$ be a projective, smooth, and geometrically connected curve over a perfect field $k$.
     Let $\mdls = \sum_{i=1}^{s} n_i P_i$ be an effective divisor where $P_1, \dots, P_s$ are distinct closed points.
     Let $U = C \setminus \mdls$ and let $d \geq \deg \mdls$.

     Suppose $\ndls_1, \dots, \ndls_l$ are pairwise coprime submoduli of $\mdls$ such that $\mdls = \sum_{j=1}^l \ndls_j$.
     Put $d_j=\deg\ndls_j$ for $1\leq j\leq l$ and $d_0=d-\deg\mdls$.
     Consider the summation morphism:
     \[
     \pi : Y:=C^{(d_1)} \times_k \dots \times_k C^{(d_l)} \times_k C^{(d_0)} \longrightarrow C^{(d)}
     \]
     defined by $(D_1, \dots, D_l, D_0) \mapsto \sum_{j=1}^l D_j + D_0$.

     Then $\pi$ is finite flat, and it is étale at the generic point of the closed subvariety
     \[
     V = \{\ndls_1\} \times_k \dots \times_k \{\ndls_l\} \times_k C^{(d_0)}\subset Y.
     \]
\end{lemma}

\begin{proof}
We may assume that $k$ is algebraically closed. The morphism $\pi$ is
finite: it is induced by the finite quotient morphism
\[
C^d/\bigl(S_{d_1}\times\cdots\times S_{d_l}\times S_{d_0}\bigr)
  \longrightarrow C^d/S_d.
\]
Its source and target are smooth $k$-schemes of pure dimension $d$.
Hence $\pi$ is flat by miracle flatness, and therefore finite flat.

Since the étale locus of $\pi$ is open and $V$ is integral, it is enough to
find one point of $V$ at which $\pi$ is étale. Choose an effective
divisor $D_0\in U^{(d_0)}$, and set
\[
v=\bigl(\ndls_1,\ldots,\ndls_l,D_0\bigr)\in V,
\qquad Q=\pi(v)=\mdls+D_0.
\]
The summands $\ndls_1,\ldots,\ndls_l,D_0$ have pairwise disjoint
supports, since the $\ndls_j$ are pairwise coprime and
$D_0$ is supported in $U$. Consider the quotient maps
\[
C^d\xrightarrow{\ p\ }
Y
\xrightarrow{\ \pi\ }C^{(d)},
\]
where $p$ groups the ordered points into blocks, and put $q=\pi\circ p$.
The same smoothness and miracle-flatness argument shows that $p$ and $q$
are finite flat, of degrees $h:=\prod_{j=0}^{l}d_j!$ and $d!$, respectively.
Write
\[
Q=\sum_{P\in\operatorname{Supp}(Q)}m_P P,
\qquad \mu:=\prod_{P\in\operatorname{Supp}(Q)}m_P!.
\]

The fiber $q^{-1}(Q)$ has $d!/\mu$ distinct points, corresponding to the
ordered lists representing $Q$. Its coordinate algebra $R$ has
$k$-dimension $d!$, since $q$ is finite flat of degree $d!$. Being
finite-dimensional over $k$, $R$ is Artinian. By \stackstag{00JB}, it
decomposes into its localizations at maximal ideals, which correspond to
the points of the fiber:
\[
R=\prod_{x\in q^{-1}(Q)}R_x,
\qquad R_x=\mathcal O_{q^{-1}(Q),x}.
\]
The action of $S_d$ is transitive on these points and induces isomorphisms
between the factors $R_x$. Therefore each has dimension
\[
\dim_k R_x=\frac{d!}{d!/\mu}=\mu.
\]

Let $A=\mathcal O_{\pi^{-1}(Q),v}$, so that $\Spec(A)$ is the local
part of the fiber at $v$, and let $B$ be the coordinate algebra of
$C^d\times_Y\Spec(A)$. As $p$ is finite flat of degree $h$, $B$ is a free
$A$-module of rank $h$. Hence
\[
\dim_k B=h\dim_k A.
\]
On the other hand, since the summands have disjoint supports, each point
$P$ of $Q$ occurs in exactly one block, with multiplicity $m_P$. The block
corresponding to $\ndls_j$ therefore has
$d_j!/\prod_{P\in\operatorname{Supp}(\ndls_j)}m_P!$ distinct orderings,
and the analogous formula holds for the block corresponding to $D_0$.
Multiplying these counts over all blocks gives exactly $h/\mu$ ordered
lists above $v$. Thus
\[
B\cong\prod_{\substack{x\in q^{-1}(Q)\\p(x)=v}}R_x,
\qquad
\dim_k B=\frac{h}{\mu}\,\mu=h.
\]
It follows that $\dim_k A=1$, and therefore $A=k$. The fiber of $\pi$ is
thus étale at $v$. Since $\pi$ is flat and of finite presentation, the
fiber criterion for étaleness (\stackstag{0476}) implies that $\pi$ is
étale at $v$, as required.
\end{proof}

\begin{cor}\label{cor:etale_at_generic_point}
For $d\geq\deg\mdls$, the addition morphism
\[
+\colon C^{(\deg\mdls)}\times_k C^{(d-\deg\mdls)}\longrightarrow C^{(d)}
\]
is finite flat, and it is \`etale at the generic point of
\[
\{\mdls\}\times_k C^{(d-\deg\mdls)}
   \subset C^{(\deg\mdls)}\times_k C^{(d-\deg\mdls)}.
\]
\end{cor}
\begin{proof}
This is the case $l=1$ and $\ndls_1=\mdls$ of \Cref{lemma:takeuchi_lemma}.
\end{proof}

We now pass from tameness on $X_{\mdls}$ to tameness on $X_{\mdls,d}$.

\begin{cor}\label{cor:tame-ramification-symmetric-product}
Let $\mdls$ be a modulus, let $d\geq\deg\mdls$, and let $\cP$ be a
$G$-torsor on $U$ whose ramification is bounded by $\mdls$. Assume that
$\cP^{[\deg\mdls]}$ is tamely ramified at the generic point
$\eta_{\mdls}$ of $E_{\mdls}$. Then $\cP^{[d]}$ is tamely ramified at
the generic point $\eta_0=\eta_{\mdls,d}$ of the exceptional divisor
$E_0=E_{\mdls,d}$ over $Z_0(\mdls,d)$.
\end{cor}
\begin{proof}
Work on $X_{\mdls}\times_k U^{(d-\deg\mdls)}$, with projections $p_1,p_2$,
and let $\theta$ be the generic point of
$E_{\mdls}\times_k U^{(d-\deg\mdls)}$. Stability of tame ramification
under base change (\stackstag{0EYD}) shows that
$p_1^{-1}\cP^{[\deg\mdls]}$ is tamely ramified at $\theta$. The torsor
$p_2^{-1}\cP^{[d-\deg\mdls]}$ is defined over the whole product, so it is
unramified, and hence tamely ramified, at $\theta$. By
\Cref{lemma:bounding_ramification_of_contracted_product},
\[
p_1^{-1}\cP^{[\deg\mdls]}
  \otimes^G p_2^{-1}\cP^{[d-\deg\mdls]}
\]
is therefore tamely ramified at $\theta$.

Consider the addition morphism
\[
+\colon C^{(\deg\mdls)}\times_k U^{(d-\deg\mdls)}\longrightarrow C^{(d)}.
\]
Its scheme-theoretic inverse image of $Z_0(\mdls,d)$ is
$\{\mdls\}\times_k U^{(d-\deg\mdls)}$. Flat base change for blowups
(\Cref{theorem:blowup_flat_base_change}) therefore gives the Cartesian
square
\[
\begin{tikzcd}[column sep=large,row sep=large]
X_{\mdls}\times_k U^{(d-\deg\mdls)}
  \arrow[r,"\widetilde{+}"] \arrow[d]
  & X_{\mdls,d} \arrow[d] \\
C^{(\deg\mdls)}\times_k U^{(d-\deg\mdls)}
  \arrow[r,"+"]
  & C^{(d)}.
\end{tikzcd}
\]
The vertical arrows are the blowup maps.
The induced morphism
$E_{\mdls}\times_k U^{(d-\deg\mdls)}\to E_0$ is flat, being a base
change of $+$. By going down for flat morphisms (\stackstag{00HS}), it
maps the generic point $\theta$ to the generic point $\eta_0$.
The left vertical map sends $\theta$ to the generic point of
$\{\mdls\}\times_k U^{(d-\deg\mdls)}$, where addition is étale by
\Cref{cor:etale_at_generic_point}. Since $\widetilde{+}$ is its base
change, it is étale at $\theta$.

On the dense open $U^{(\deg\mdls)}\times_k U^{(d-\deg\mdls)}$,
$\widetilde{+}$ agrees with addition. Hence
\Cref{prop:symmetric_powers_on_curves} gives
\[
(\widetilde{+})^{-1}\cP^{[d]}\cong
p_1^{-1}\cP^{[\deg\mdls]}
  \otimes^G p_2^{-1}\cP^{[d-\deg\mdls]}.
\]
This torsor is tamely ramified at $\theta$. Tameness therefore descends
along $\widetilde{+}$ to $\eta_0$ by
\Cref{lemma:descend_ramification_along_etale}.
\end{proof}


\begin{proof}[Proof of \Cref{theorem:SymmetricPowerOfSheafIsTamelyRamified}]
By flat base change, \Cref{lemma:weakly_unramified_base_change}, and
\Cref{prop:equivalence_ramification_bounded_by_modulus}, we may assume
that $k$ is algebraically closed.

For each single-point summand $\ndls=n_P P$ of $\mdls$,
\Cref{theorem:torsors_after_blowup_is_tame} gives tameness of
$\cP^{[\deg\ndls]}$ at $\eta_\ndls$. Applying
\Cref{cor:moduli-reduction} repeatedly gives tameness of
$\cP^{[\deg\mdls]}$ at $\eta_\mdls$.
\Cref{cor:tame-ramification-symmetric-product} then gives tameness of
$\cP^{[d]}$ at $\eta_0$ for every sufficiently large $d$. Since $H$ is a
nonempty open subset of $E_0$, it has generic point $\eta_0$, proving
the theorem over $\bar{k}$. The initial base-change reduction descends this
tameness to the original field $k$.
\end{proof}

